% BOOK-FILE-ID: f-d55c003e2829
% section: 漸化式の分類

\section{解法の分類}\label{節：解法の分類}

ここまで様々な漸化式を扱ってきた.最後に,それらを振り返って整理しておこう.

まず,本書で扱った解法は大きく2つに分けられる.
1つは,式変形によって既知の形をつくる解法である.
4つの基本パターンを行き先として,応用パターン,発展的な解法へと広げてきた.
もう1つは,式そのものが持つ構造を使う解法である.
対称式や三角関数の恒等式,変わらない量,初項の選び方などを手掛かりにした.
ガウス記号を含む式のように,両者が重なることもある.
\BookFigClassification

前者を振り返ると,どの解法も同じ流れをたどっていたことに気付くだろう.
うまく置き換えると,難しい操作が足し算や掛け算に変わる.そこで解いて,元に戻す.
置き換える相手は不動点からのずれ,逆数や対数,2つの項の組み合わせとさまざまだったが,多くの場合,行き先は基本パターンであった.
この共通の流れがなぜ現れるのかは,付録の\sectionref*{節：本書の解法を見渡す}で考える.

次に,「漸化式を解く」とは何を求めることなのかを考えよう.
多くの場合は\bookterm{general}{一般項}を簡単な式で表すことが目標だが,それだけではない.
一般項を和や積で表すところまでしかできないこともあるし,周期や止まる値,整数になるかどうかといった性質を調べるのが目標のこともある.

では,「簡単な式で表す」とはどういうことだろうか.
本書では,高校数学で扱う関数を使い,\,$n$が大きくなっても計算の回数が増えない式で一般項を表すことを目標とする.\footnote{本書ではガウス記号や場合分けも使う.どのような式を答えとして認めるかは,使ってよい関数の範囲を決めて初めて考えられる.}
注意したいのは,本書の方法でそのような式が見つからなかったからといって,他の方法でも見つからないとは言えないことである.

最後に,ここで示した分類はあくまで解法を探すための案内であり,これに当てはめれば必ず解けるというものではない.
そのような式が見つからなくても数列は定義されているし,\bookterm{limit}{極限}や整数の性質を調べることはできる.
\bookproblemref{practice-entrance}{6}{発展問6}では,一般項を求めずに数列の減り方を調べる.
値の変化を分析することや,解けない漸化式についてじっくりと考えること,それもまた漸化式の面白さの1つなのである.
一般項が見つからない場合の考え方は,付録の\sectionref*{節：置き換えられないとき}でも扱う.
